\subsection{2. Ein unveränderter
Modellvertrag}\label{ein-unveruxe4nderter-modellvertrag}

Es werden keine Hopping-, Massen-, Ladungs- oder Projektionsglieder zu H
hinzugefügt:

\[\adjustbox{max width=\linewidth}{$\displaystyle 
H=\Delta N_b+g(Q_++Q_-),\quad Q_+=\sum_A b_A^\dagger P_A,
\quad Q_-=Q_+^\dagger,
$}\] \[\adjustbox{max width=\linewidth}{$\displaystyle 
P_A=\sum_{i<j}W_{A,ij}f_jf_i,\quad N=N_f+2N_b,
\quad WW^\dagger=8I_{60}.
$}\]

Es gibt 64 CAR-Fermionmoden, 60 CCR-Bosonmoden und 480 von null
verschiedene reelle W-Einträge mit ihren ursprünglichen Vorzeichen.
\(\Delta>0\), g ist reell; die Zahlen beziehen sich auf
\(g/\Delta=1/20\). Die Tensorquelldatei ist durch SHA-256
\codeword{3f00a0892157a691e4ef26920c3702772c7bf76a250fad98ff04a079bb64b763}
fixiert.

Die Fockrealisierung, H, der sektorenübergreifende Energievergleich und
der Prüfparameter sind weiterhin Modellvoraussetzungen. Die geometrische
oder arithmetische Herkunft eines W-Tensors leitet diese Voraussetzungen
nicht allein her. Insbesondere sind die fünf elementaren CAR in einer
Spinor-Matrixkonstruktion nicht mit den hier verwendeten 64 Fockmoden
gleichzusetzen.

\subsection{3. Was tatsächlich erneut geprüft
wurde}\label{was-tatsuxe4chlich-erneut-gepruxfcft-wurde}

\subsubsection{3.1 Grundzustand: frische Berechnung statt bloßer
Statusübernahme}\label{grundzustand-frische-berechnung-statt-blouxdfer-statusuxfcbernahme}

Die ursprünglichen Programme wurden mit fixierten Dateihashes in eine
eigene Arbeitskopie übernommen. Die Paar-Konfigurationen wurden frisch
mit neu kompiliertem Quellcode berechnet:

{\def\LTcaptype{none} % do not increment counter
\begin{longtable}[]{@{}
  >{\raggedright\arraybackslash}p{(\linewidth - 4\tabcolsep) * \real{0.2727}}
  >{\raggedleft\arraybackslash}p{(\linewidth - 4\tabcolsep) * \real{0.3636}}
  >{\raggedleft\arraybackslash}p{(\linewidth - 4\tabcolsep) * \real{0.3636}}@{}}
\toprule\noalign{}
\begin{minipage}[b]{\linewidth}\raggedright
Ordnung
\end{minipage} & \begin{minipage}[b]{\linewidth}\raggedleft
Vollständig enumerierte Bosonkonfigurationen
\end{minipage} & \begin{minipage}[b]{\linewidth}\raggedleft
Normquadrat von \(Q_+^kF\)
\end{minipage} \\
\midrule\noalign{}
\endhead
\bottomrule\noalign{}
\endlastfoot
0 & analytischer Ausgangszustand & 1 \\
1 & 60 Kanäle / 480 Paare & 480 \\
2 & 1830 & 439680 \\
3 & 37820 & 575078400 \\
4 & 595665 & 952296652800 \\
\end{longtable}
}

Zusätzlich wurden unabhängige Wortentwicklungen, Überlaufgrenzen, die
vollständige Zweikörper-Casimiridentität, rationale Sektorvergleiche und
die unabhängige Symmetriespurrechnung wiederholt. Die vier ausgewählten
ursprünglichen Zertifikate bestehen normal und unter \codeword{-OO} mit
identischen JSON-Bytes. Dies ist eine gezielte erneute
Grundzustandsprüfung, nicht die Behauptung, alle historischen
Repository-Tests erneut ausgeführt zu haben. Der unveränderte alte
Normprüfer gibt unter dem aktuellen NumPy eine ComplexWarning bei der
Ganzzahlkonversion aus. Die neuen Prüfer bestätigen vor jeder Konversion
exakt verschwindende Imaginärteile und ganzzahlige Realteile des
gepinnten W. Die Warnung ist gespeichert und wird nicht als verlorener
Tensoranteil oder als unterdrückter Prüffehler ausgegeben.

Der modellinterne Satz bleibt bestehen: Für \(0<|g|/\Delta\le1/20\) ist
der globale Grundzustand \(\Omega\) eindeutig, hat N=64 und ist
Spin(10)×SU(4)-invariant. Am Zahlenprüfpunkt:

\[\adjustbox{max width=\linewidth}{$\displaystyle 
-1.158089\Delta<E_0<-1.129636\Delta,
\qquad \operatorname{gap}(H)>0.007737\Delta.
$}\]

Der volle Grundzustandsvektor ist damit nicht ausgerechnet. Die fünf
Versuchsvektoren sind keine behauptete invariante Fünferbasis für H.

\subsubsection{3.2 Das neu eingegangene
Polergebnis}\label{das-neu-eingegangene-polergebnis}

Der vollständige technische Bericht v1.6.3 und sein 248-zeiliger Prüfer
wurden gelesen; Quellen, Bericht, Ergebnismatrix und Prüfpaket wurden
unverändert archiviert. Der Prüfer wurde normal und optimiert
wiederholt: \textbf{317 Prüfbedingungen, identische JSON-Bytes auch zum
gelieferten Bericht}.

Der Herkunftspin des Prüfers ist
\codeword{fdbcabd244450c302182086d67c68284634c994e3fee026200c42c508f731654}.
Er verwendet denselben ursprünglichen
\codeword{verify\_hole.py}-Quellstand wie die bisherige Fortsetzung. Die
neuen eigenen Sektorvergleiche wurden zusätzlich unabhängig mit
rationalen LDL-Pivots berechnet.

Die Prüfzahlen zählen auch Komponenten und Wiederholungen; sie sind
keine Zahl unabhängiger Entdeckungen. Die analytischen Argumente sind
ausgeschrieben, aber nicht vollständig in einem Beweisassistenten
formalisiert.

\subsection{4. Der einfache native Anschluss: Bewegungsgleichung statt
Umdeutung}\label{der-einfache-native-anschluss-bewegungsgleichung-statt-umdeutung}

Erweitere jede W-Zeile zur antisymmetrischen Matrix \(M_A\) mit
\((M_A)_{ij}=W_{A,ij}\) für i\textless j. Definiere

\[\adjustbox{max width=\linewidth}{$\displaystyle 
D_r=\sum_{A,j}(M_A)_{rj}b_Af_j^\dagger,
\qquad \phi_r=D_r/\sqrt{15}.
$}\]

Direkte CAR/CCR-Normalordnung liefert \textbf{Operatoridentitäten auf
dem endlichen Teilchenkern}, nicht nur Gleichheiten auf Testzuständen:

\[\adjustbox{max width=\linewidth}{$\displaystyle 
\boxed{[H,f_r]=-gD_r,\qquad [H,f_r^\dagger]=gD_r^\dagger,}
$}\] \[\adjustbox{max width=\linewidth}{$\displaystyle 
[N,D_r]=-D_r,\qquad \{f_r,D_s^\dagger\}=0.
$}\]

Der neue Vergleichsoperator hat also dieselbe Ladung −1 wie f. Er ist
nicht das früher auf dem leeren Referenzzustand betrachtete
\(\chi^\dagger\sim b^\dagger f^\dagger\) mit Ladung +3. Auf dem
Grundzustand liegen \(f\Omega\) und \(D\Omega\) in N=63,
\(\chi^\dagger\Omega\) dagegen in N=67. Neutrale Zeitentwicklung hebt
diese Unterscheidung nicht auf.

Die Normalordnungsprüfung erfasst alle 64 ursprünglichen f-Operatoren
und die tatsächlichen W-Vorzeichen. Für die Kompositnorm gilt exakt

\[\adjustbox{max width=\linewidth}{$\displaystyle 
\{D_r,D_s^\dagger\}=
\sum_{A,B,j,k}(M_A)_{rj}(M_B)_{sk}
\left(\delta_{AB}f_j^\dagger f_k+\delta_{jk}b_B^\dagger b_A\right).
$}\]

Mit \(\bar b=\langle N_b\rangle\) und der Grundzustandssymmetrie:

\[\adjustbox{max width=\linewidth}{$\displaystyle 
\langle\{D_r,D_s^\dagger\}\rangle=\delta_{rs}S,
\qquad S=15-\frac7{32}\bar b.
$}\]

\(\phi\) hat entsprechend Norm \(1-7\bar b/480\), nicht globale
kanonische CAR. Der normierte Zustandserwartungswert ersetzt keine
Operatorrelation.

\subsection{5. Dieselbe Antwort auf demselben
Grundzustand}\label{dieselbe-antwort-auf-demselben-grundzustand}

Für Im z\textgreater0 und \(H_n=H|_{N=n}\):

\[\adjustbox{max width=\linewidth}{$\displaystyle 
G_{rs}(z)=\langle\Omega|f_r(z+E_0-H_{65})^{-1}f_s^\dagger|\Omega\rangle
+\langle\Omega|f_s^\dagger(z-E_0+H_{63})^{-1}f_r|\Omega\rangle.
$}\]

Die innere Symmetrie macht G diagonal und alle Diagonalelemente gleich.
Schreibe mit positiven Entnahme- und Additionsmaßen auf \(\epsilon>0\)

\[\adjustbox{max width=\linewidth}{$\displaystyle 
G(z)=\int\frac{d\nu_+(\epsilon)}{z-\epsilon}
+\int\frac{d\nu_-(\epsilon)}{z+\epsilon}.
$}\]

Dann gelten exakt

\[\adjustbox{max width=\linewidth}{$\displaystyle 
Z_+=\bar b/32,\quad Z_-=1-\bar b/32,
\quad a:=\int\epsilon\,d\nu_+=\int\epsilon\,d\nu_-
=\frac{\Delta\bar b-E_0}{64}.
$}\]

Für das gesamte signierte Spektralmaß:

\[\adjustbox{max width=\linewidth}{$\displaystyle 
\boxed{m_0=1,\quad m_1=0,\quad m_2=g^2S,\quad
m_3=g^2(\Delta S+7a).}
$}\]

Die ersten beiden Momente und die Kanalgewichte stimmen mit der
eingegangenen unabhängigen Rechnung überein. \textbf{Das dritte Moment
ist der zusätzliche eigene Schritt dieser Revision.}

\subsubsection{5.1 Herleitung des dritten
Moments}\label{herleitung-des-dritten-moments}

Setze \(\mathcal L A=[A,H]\). Stationarität macht diesen Operator
symmetrisch in der positiven, nach Nullvektoren quotientierten Metrik
\((A,B)=\langle\{A^\dagger,B\}\rangle\). Direkte Normalordnung ergibt

\[\adjustbox{max width=\linewidth}{$\displaystyle 
\sum_r\{[D_r,X],D_r^\dagger\}=-14Q_+.
$}\]

Die beiden Beiträge sind \(16Q_+\) und \(-30Q_+\). Ihre Koeffizienten
folgen aus den vollständig geprüften Kontraktionen
\(\sum_{rj}M_{A,rj}M_{B,rj}=16\delta_{AB}\) und
\(\sum_{Ar}M_{A,rj}M_{A,rk}=15\delta_{jk}\). Mit \([D,N_b]=D\) und
\(\langle Q_+\rangle=(E_0-\Delta\bar b)/(2g)\) folgt die Formel.

Als unabhängige Kontrolle wurden sämtliche Formeln einschließlich m3 an
der früher exakt geschlossenen N=4/N=5-Antwort symbolisch geprüft. Diese
Kontrolle verwendet den früheren Referenzzustand nur als Test der
Identitäten, nicht als Ersatz für den nativen Grundzustand.

Die allgemeine Methode, Spektralmomente aus Bewegungsgleichungen zu
gewinnen, ist etabliert; siehe
\href{https://arxiv.org/abs/0907.1284}{Freericks und Turkowski, Phys.
Rev.~B 80, 115119}. Die hier angegebenen W-Kontraktionen und Konstanten
wurden eigenständig für dieses Modell berechnet; die zitierte Arbeit
beweist keine TFPT-Aussage.

\subsubsection{5.2 Neue engere Dichte- und
Gewichtsgrenzen}\label{neue-engere-dichte--und-gewichtsgrenzen}

Positivität der beiden Antwort-Grammatrizen beziehungsweise zweimal
Cauchy-Schwarz liefern

\[\adjustbox{max width=\linewidth}{$\displaystyle 
m_2\ge a^2\left(\frac1{Z_-}+\frac1{Z_+}\right),
$}\] \[\adjustbox{max width=\linewidth}{$\displaystyle 
\boxed{(\Delta\bar b-E_0)^2\le
60g^2\bar b(32-\bar b)(1-7\bar b/480).}
$}\]

Mit der bereits bewiesenen Energieobergrenze und \(g/\Delta=1/20\) muss
das folgende rationale Polynom positiv sein:

\[\adjustbox{max width=\linewidth}{$\displaystyle 
P(b)=\frac7{3200}b^3-\frac{61}{50}b^2
+\frac{317591}{125000}b-\frac{79754843281}{62500000000}>0.
$}\]

Seine beiden im alten zulässigen Bereich liegenden Nullstellen liegen
bei ungefähr 0.842846697 und 1.245655664. Exakte rationale
Wurzelisolation ergibt die nach außen gerundete strenge Schranke

\[\adjustbox{max width=\linewidth}{$\displaystyle 
\boxed{0.842846<\bar b<1.245656.}
$}\]

Damit:

{\def\LTcaptype{none} % do not increment counter
\begin{longtable}[]{@{}
  >{\raggedright\arraybackslash}p{(\linewidth - 2\tabcolsep) * \real{0.4286}}
  >{\raggedleft\arraybackslash}p{(\linewidth - 2\tabcolsep) * \real{0.5714}}@{}}
\toprule\noalign{}
\begin{minipage}[b]{\linewidth}\raggedright
Größe pro Mode
\end{minipage} & \begin{minipage}[b]{\linewidth}\raggedleft
Strenges offenes Intervall
\end{minipage} \\
\midrule\noalign{}
\endhead
\bottomrule\noalign{}
\endlastfoot
Gesamtes Additionsgewicht & (0.0263389375, 0.03892675) \\
Gesamtes Entnahmegewicht & (0.96107325, 0.9736610625) \\
\(\langle\{\phi,\phi^\dagger\}\rangle\) & (0.981834183333\ldots,
0.987708495833\ldots) \\
Frühere \(\chi\)-Kompositnorm, nicht \(\phi\) & (0.040386370833\ldots,
0.059687683333\ldots) \\
\end{longtable}
}

Das sind Erwartungswerte und integrierte Spektralgewichte, keine
direkten Nachweise von Produktionsraten oder von bereits verfügbaren
Messinstrumenten.

\subsection{6. Der Polsatz und seine zusätzliche
Verschärfung}\label{der-polsatz-und-seine-zusuxe4tzliche-verschuxe4rfung}

Der eingegangene Beweis verwendet \(u_{k,r}=Q_+^kf_rF=f_rQ_+^kF\).
Invarianz und Besetzung ergeben

\[\adjustbox{max width=\linewidth}{$\displaystyle 
\langle u_{k,r},u_{k,s}\rangle=
\delta_{rs}\frac{64-2k}{64}\|Q_+^kF\|^2.
$}\]

Die Normen lauten 1, 465, 412200, 521164800, 833259571200. Die daraus
gebildete fünfdimensionale Variationsmatrix gibt 64 unabhängige
Richtungen unter \(-1.095812\Delta\). Das gesamte N=63-Komplement des
Nullbosonraums liegt über \(-3\Delta/4\). Minimax begrenzt den niedrigen
Raum auf genau 64 Dimensionen; seine injektive symmetrieverträgliche
Projektion auf die irreduzible duale 64 erzwingt ein einziges
Energieniveau.

Das beweist Existenz und Isolation. Seine Sichtbarkeit folgt aus den
positiven Entnahmemomenten. Mit

\[\adjustbox{max width=\linewidth}{$\displaystyle 
d=0.007737\Delta,\quad c=0.379636\Delta,
\quad a_{\max}=\frac{1.245656+1.158089}{64}\Delta
$}\]

gilt

\[\adjustbox{max width=\linewidth}{$\displaystyle 
Z_{\rm low}>\frac{c(1-1.245656/32)-a_{\max}}{c-d}
=\frac{40912436089}{46487375000}
=0.8800762806891118\ldots.
$}\]

Außerdem ist der niedrige Pol die kleinste Entnahmeenergie. Daher
\(a\ge\epsilon_{\rm low}Z_-\), also bereits ohne vollständige
Polauswertung

\[\adjustbox{max width=\linewidth}{$\displaystyle 
\frac{\epsilon_{\rm low}}\Delta
<\frac{1.245656+1.158089}{64-2(1.245656)}
=\frac{2403745}{61508688}=0.039079763821332\ldots.
$}\]

Diese Verschärfungen entstehen aus dem \textbf{Zusammenführen
kompatibler Beweise}, nicht aus einem neuen gewählten Parameter. Das
Restgewicht des gesamten CAR-Maßes ist somit kleiner als
0.119923719311\ldots{} .

\subsection{7. Einfache Organisation, aber keine falsche
Zwei-Linien-Lösung}\label{einfache-organisation-aber-keine-falsche-zwei-linien-luxf6sung}

Die ersten zwei orthonormalen Operatoren sind f und \(D/\sqrt S\). Die
erste Resolventenreduktion hat deshalb die exakte Form

\[\adjustbox{max width=\linewidth}{$\displaystyle 
\boxed{G(z)=\frac1{z-\displaystyle\frac{g^2S}{z-a_1-\Sigma_2(z)}}},
\qquad a_1=\Delta+\frac{7a}{S}.
$}\]

\(\Sigma_2\) ist die Resolvente des verbleibenden nativen Operatorraums,
gekoppelt an den dazu orthogonalen Rest von \(\mathcal L(D/\sqrt S)\).
Es wurde kein äußeres Bad hinzugefügt. Dies ist eine genaue Ordnung der
Rechnung, \textbf{keine abgeschlossene Berechnung von \(\Sigma_2\)} und
keine Vereinigung sämtlicher TOE-Aufgaben in einer einzigen bereits
gelösten Funktion.

\subsubsection{7.1 Warum man die Rückwirkung nicht exakt weglassen
darf}\label{warum-man-die-ruxfcckwirkung-nicht-exakt-weglassen-darf}

Angenommen, es gäbe nur je eine Entnahme- und Additionslinie mit
Energien \(\epsilon_-,\epsilon_+\). Symmetrie macht diese für alle r
gleich. Dann liefern die Bewegungsgleichungen auf demselben Grundzustand

\[\adjustbox{max width=\linewidth}{$\displaystyle 
D_r\Omega=-\frac{\epsilon_-}{g}f_r\Omega,
\qquad D_r^\dagger\Omega=\frac{\epsilon_+}{g}f_r^\dagger\Omega.
$}\]

Summe nach Multiplikation mit \(f_r^\dagger\) beziehungsweise \(f_r\)
und N=64 ergeben

\[\adjustbox{max width=\linewidth}{$\displaystyle 
H\Omega=\left[-32\epsilon_-+
(\Delta-\epsilon_++\epsilon_-)N_b\right]\Omega.
$}\]

Ein Eigenzustand mit negativer Energie kann keine feste Bosonenzahl
haben: Dann wäre \(\langle Q_++Q_-\rangle=0\) und seine Energie
\(\Delta\langle N_b\rangle\ge0\). Daher sind \(\Omega\) und
\(N_b\Omega\) unabhängig. Die angenommene Zweilinienform erzwingt
\(\epsilon_+-\epsilon_-=\Delta\).

Andererseits geben \(m_0=1,m_1=0\) für ein Zweilinienmaß
\(m_3/m_2=\epsilon_+-\epsilon_-\). Die exakt bestimmte Formel lautet
jedoch

\[\adjustbox{max width=\linewidth}{$\displaystyle 
\frac{m_3}{m_2}=\Delta+7a/S>\Delta.
$}\]

Widerspruch. Damit ist \(\Sigma_2\not\equiv0\) bewiesen. Mindestens ein
Kanal besitzt mehr als eine Energielinie. Der dominante isolierte
Entnahmepol und diese unvermeidliche Reststruktur widersprechen einander
nicht. Eine Zweilinienform könnte höchstens eine zu zertifizierende
Näherung sein.

